\section{The Orbit-Specificity Boundary}\label{sec:barrier}

This section establishes that the orbit-specificity barrier is the sole
remaining obstacle to FF-MC, and that it is structurally identical to the
barrier for Mullin's Conjecture over $\ZZ$. We present three formally
verified dead ends that delineate the boundary of current techniques.

\subsection{Population versus orbit equidistribution}

The central dichotomy is between two types of equidistribution:

\begin{definition}[Population equidistribution]
Among all monic irreducible polynomials of degree $d$, the residues
modulo a fixed irreducible $Q$ are equidistributed across the
$|\GF{p}{\deg Q}^\times|$ residue classes, with error
$O(p^{d/2})$ from the Weil bound.
\end{definition}

\begin{definition}[Orbit equidistribution]
The specific walk $\ffprod(n) \bmod Q$ visits each residue class in
$\GF{p}{\deg Q}^\times$ with asymptotic frequency $1/(p^{\deg Q} - 1)$.
This is the content of the function field DSL (FF-DSL).
\end{definition}

Population equidistribution is unconditional over $\Fpt$ (from the Weil
bound, which is a theorem of Weil~\cite{Weil48}). Orbit equidistribution
remains open. The gap between them is precisely the orbit-specificity
barrier.

\subsection{Dead End \#127: Weil bound insufficient}

\begin{theorem}[Weil Does Not Close the Gap]\label{thm:de127}
\leancheck{dead\_end\_127\_weil\_not\_sufficient}
The Weil bound implies population equidistribution but does not imply
FF Dynamical Hitting ($\FFDH$). The gap persists because:
\begin{enumerate}[label=(\roman*)]
  \item The Weil bound controls sums of the form
    $\sum_{\deg f = d, f \text{ monic}} \chi(f \bmod Q)$,
    which average over \emph{all} monic polynomials of degree $d$.
  \item $\FFDH$ requires control of
    $\sum_{n < N} \chi(\ffprod(n) \bmod Q)$,
    which is a sum along a \emph{specific deterministic orbit}.
  \item The polynomials $\ffprod(n) + 1$ are products of specific
    irreducibles, not random draws from the population.
\end{enumerate}
\end{theorem}

The structural reason is that $\ffprod(n) + 1$ has degree
$\sum_{k \leq n} \deg(\ffseq(k))$, which grows rapidly. Among monic
polynomials of this degree, $\ffprod(n) + 1$ is essentially the
\emph{unique representative} of its congruence class---exactly as
$P(n) + 1$ is over $\ZZ$.

\subsection{Dead End \#129: Large monodromy structurally false}

The monodromy approach attempts to bypass the Weil bound limitation by
using non-abelian Galois theory. The Galois group
$\mathrm{Gal}(\mathrm{Split}(\ffprod(n)+1)/\Fp)$ acts on the roots of
$\ffprod(n)+1$, and its cycle structure determines the degrees of
irreducible factors.

\begin{definition}[FF Large Monodromy (FFLM)]
\leancheck{FFLargeMonodromy}
The Galois group of $\ffprod(n)+1$ eventually contains the alternating
group: $|\mathrm{Gal}| \geq (\deg \ffprod(n))!/2$ for all sufficiently
large $n$.
\end{definition}

If FFLM held, the Deligne equidistribution theorem would give orbit-specific
information about the factorization pattern of $\ffprod(n) + 1$, going
beyond what the Weil bound provides.

\begin{theorem}[FFLM is Structurally False]\label{thm:de129}
\leancheck{dead\_end\_129\_monodromy\_dead}
The greedy FF-EM construction produces products dividing $t^{p^d} - t$
for some $d$ (since the product of all monic irreducibles of degree
dividing $d$ equals $t^{p^d} - t$). Adding $1$ gives a polynomial whose
splitting field is contained in $\GF{p}{N}$ for some $N$, hence the Galois
group is cyclic: $\mathrm{Gal} \cong \ZZ/N\ZZ$, generated by Frobenius.
In particular, the Galois group is \emph{abelian}, not $S_d$ or $A_d$.
\end{theorem}

The consequence is that the monodromy approach reduces to abelian Galois
theory, which is captured by Dirichlet characters, which is controlled by
the Weil bound, which gives population equidistribution---the same
information we already have. The monodromy route is circular.

\subsection{Dead End \#130: Generation does not imply coverage}

A natural hope is that the subgroup escape property (the walk steps
generate the full group) might imply dynamical hitting (the walk visits
every element). Over both $\ZZ$ and $\Fpt$, subgroup escape is proved
(PRE over $\ZZ$, Weil over $\Fpt$). The following counterexample shows
this is insufficient.

\begin{theorem}[Generation $\neq$ Coverage]\label{thm:de130}
\leancheck{dead\_end\_130\_generation\_ne\_coverage}
Consider $\ZZ/4\ZZ$ with the alternating walk using steps $\{1, 3\}$.
\begin{enumerate}[label=(\roman*)]
  \item The set $\{1, 3\}$ generates $\ZZ/4\ZZ$ as an additive group.
  \item The alternating walk $0, 1, 0, 1, \ldots$ (adding $1$, then $3$,
    then $1$, etc.) visits only $\{0, 1\}$.
  \item The elements $2$ and $3$ are never reached.
\end{enumerate}
\end{theorem}

This demonstrates that even when the walk steps can reach everywhere (in the
algebraic sense of group generation), the \emph{deterministic accumulation}
of steps can miss elements. The gap between generation (SE) and coverage
(DH) is structural, not an artifact of the specific group.

\subsection{The orbit barrier thesis}

\begin{theorem}[Orbit Barrier Thesis]\label{thm:barrier-thesis}
\leancheck{orbit\_barrier\_thesis}
The orbit-specificity barrier has the following properties:
\begin{enumerate}[label=(\roman*)]
  \item All structural prerequisites for FF-MC are proved unconditionally:
    injectivity, degree growth, finiteness of irreducibles per degree,
    coprimality cascade.
  \item $\FFDH$ alone implies $\FFMC$.
  \item $\FFDH$ is the sole remaining hypothesis.
  \item The barrier is not solvable by population statistics (Dead End \#127).
  \item The barrier is not solvable by Galois theory (Dead End \#129).
  \item The barrier is not solvable by algebraic generation (Dead End \#130).
  \item The barrier dissolves under stochastic perturbation ($\varepsilon > 0$).
  \item The barrier is sharp at $\varepsilon = 0$ (phase transition).
\end{enumerate}
\end{theorem}

\subsection{Setting independence}

The orbit barrier thesis has a striking corollary: the barrier is
\emph{setting-independent}. Over $\ZZ$, Mullin's Conjecture has two
potential obstacles:
\begin{enumerate}
  \item The analytic number theory input (PrimesEquidistributedInAP),
    which is standard but requires substantial infrastructure.
  \item The orbit-specificity barrier (DSL).
\end{enumerate}
Over $\Fpt$, obstacle~(1) is eliminated (the Weil bound makes ANT free),
but obstacle~(2) persists unchanged. This proves that the difficulty of
Mullin's Conjecture is \emph{entirely} in the orbit-specificity barrier,
and not in the analytic number theory.

\begin{theorem}[Setting Independence]\label{thm:setting-independent}
\leancheck{barrier\_setting\_independent}
\begin{enumerate}[label=(\roman*)]
  \item Over $\Fpt$: ANT is free (Weil $\Rightarrow$ PE).
  \item Over $\Fpt$: orbit-specificity remains ($\FFDH$ open, $\FFDH
    \Rightarrow \FFMC$ proved).
  \item The stochastic perturbation dissolves the barrier (FFStochasticMC
    proved).
  \item The deterministic limit does not (FFPhaseTransition proved).
\end{enumerate}
\end{theorem}

\subsection{Summary comparison with the integer case}

\begin{center}
\renewcommand{\arraystretch}{1.3}
\begin{tabular}{lcc}
\toprule
\textbf{Aspect} & \textbf{Over $\ZZ$} & \textbf{Over $\Fpt$} \\
\midrule
Population equidistribution & Open (WPNT) & \textbf{Free} (Weil) \\
Subgroup escape & Proved (PRE) & Proved (Weil) \\
Coprimality cascade & Proved & Proved \\
CME $\Rightarrow$ MC chain & Proved & Proved \\
Orbit specificity (DSL) & \textbf{Open} & \textbf{Open} \\
Walk character sums & Not standard & Not standard \\
Stochastic MC ($\varepsilon > 0$) & Proved & Proved \\
$\Phi_3$ exclusion & No analog & \textbf{Proved} \\
\bottomrule
\end{tabular}
\end{center}

The table confirms the thesis: the function field setting eliminates the
``easy'' part (ANT for PE) and adds new structural information ($\Phi_3$
exclusion), but the ``hard'' part (orbit specificity $=$ DSL) is identical.
